\documentclass[10pt,a4paper]{article}

\usepackage[margin=0.85in]{geometry}
\usepackage{amsmath,amssymb}
\usepackage{booktabs}
\usepackage{enumitem}
\usepackage{microtype}
\usepackage{xcolor}
\usepackage{listings}
\usepackage[hidelinks]{hyperref}

\setlist{nosep,leftmargin=1.2em}
\lstset{basicstyle=\small\ttfamily,frame=single,numbers=none,
        breaklines=true,framerule=0.4pt,
        backgroundcolor=\color{black!2}}

\title{\vspace{-3em}\textbf{Software defence against RAM-replay attacks on
        TME-only commodity hardware: tagged-pointer MACs in BEAM}}

\author{%
  Claude Opus 4.7[1m]\\\small{Anthropic}
  \and
  Sam Williams\\\small{\texttt{sam@arweave.org}}
}

\date{April 2026 --- design note (companion to LapEE)}

\begin{document}
\maketitle
\vspace{-1.8em}

\begin{abstract}
\noindent
Intel Total Memory Encryption (TME) provides confidentiality for DRAM contents
but not freshness: a physical-access attacker can capture ciphertext at a
given physical address and replay it later, and the CPU will decrypt the
captured value identically. The proper defence is hardware memory-integrity
trees (SEV-SNP RMP, TDX MAC trees), which are unavailable on the commodity
laptop hardware that the LapEE project targets. We describe a
software-only defence that detects RAM replay against a running BEAM virtual
machine by storing a per-pointer truncated MAC of the pointed-to bytes in
the unused high bits of every term-bearing pointer. Because BEAM terms form
a graph rooted in a small set of process-control structures, an interior
replay is detected on the next dereference of the parent pointer, and the
attacker's effective surface collapses to graph roots. Combined with
keeping process-control blocks resident in CPU cache, the residual surface
is reduced to a handful of mostly-static pages. We sketch the construction,
its bit budget, the engineering and performance trade-offs, and the
unresolved design questions, and conclude with a path to a measured
prototype.
\end{abstract}

\section{Threat model}

\textbf{Adversary.} A physical-access attacker who can interpose on the
DRAM channel of a TME-protected node --- via a memory bus tap, a malicious
DMA peripheral, an FPGA interposer between the CPU package and DIMMs, or
a sufficiently capable cold-boot extraction. The attacker reads and
writes ciphertext at physical addresses of their choosing during the
execution of a target program.

\textbf{Goal.} Influence the output of an in-flight computation on a
LapEE-attested HyperBEAM node by making it process stale data without
detection. Concretely: rewind a memory location holding a request
parameter, an authorisation token, a counter, or any term-graph node
participating in an attested workload, and have the node sign or commit
to the resulting (now-incorrect) computation as if nothing happened.

\textbf{Out of scope.} Replay against the TPM itself --- LapEE already
defends against this with per-attestation nonces signed inside the TPM2
quote. Replay against the AK private key --- the AK lives in the TPM
and is never exposed to RAM. The class of attack we are defending here
is replay against the \emph{computation} happening in BEAM userspace,
where the TPM has already certified an honest boot but a runtime memory
adversary subverts ongoing work.

\section{Why TME alone is insufficient}

TME on Intel client silicon is constructed as AES-XTS over the DRAM bus
with a per-boot key and a tweak derived from the physical address. The
construction has two properties relevant here:

\begin{itemize}[label=---]
\item \emph{Same address, same key, same ciphertext, same plaintext.} The
encryption is deterministic; ciphertext at address $P$ at time $t_1$
decrypts identically to ciphertext at $P$ at time $t_2$ as long as the
key has not rotated. The key only rotates on power cycle.
\item \emph{No integrity tag.} TME stores no MAC, no version counter,
no Merkle commitment. The CPU has no way to detect that the ciphertext
at $P$ was substituted for an older ciphertext at the same $P$.
\end{itemize}

The freshness gap is exactly what AMD SEV-SNP's reverse-map table (RMP)
and Intel TDX's memory-encryption MAC trees were designed to close. They
are not yet available on the commodity laptop class targeted by LapEE
(Framework 13 with Intel 12th/13th-gen Core silicon ships TME but not
TDX; AMD Ryzen variants ship SME but not SEV-SNP for client SKUs).

In the absence of hardware integrity, BEAM's natural memory churn ---
generational copying GC, per-request ephemeral processes, per-scheduler
allocator separation --- already imposes some practical difficulty on a
replay attacker, because the window in which a captured ciphertext
remains valid for a given logical object is bounded by the next minor GC
cycle (typically tens of milliseconds under load). We have proposed
amplifying this churn through configuration tuning. But this is a
probabilistic defence: with enough profiling and timing, a sufficiently
patient attacker can succeed. The defence we describe here is intended
to make the attack \emph{detectable} in addition to \emph{difficult}.

\section{Construction}

\subsection{Tagged pointers in BEAM}

BEAM represents Erlang terms in single machine words. The bottom 2--4
bits encode a major term tag (immediate, list cell, boxed, etc.); for
boxed and list-cell terms the remaining bits hold a pointer to a heap
allocation. On x86\_64 Linux, virtual addresses use at most 47 bits in
canonical form, so a term-bearing pointer naturally has $\sim$~16
unused high bits. On a laptop addressing $\le$~32~GB of RAM the live
address space is much smaller --- 35--36 bits of address suffice ---
which leaves $\sim$~24 high bits genuinely free.

\subsection{Per-pointer integrity tag}

We propose to use those high bits as a truncated message authentication
code (MAC) over the bytes that the pointer points at:

\begin{equation*}
\mathit{tag}(p) = \mathrm{MAC}_K(\mathit{content}(p))[0..n)
\end{equation*}

where $K$ is a per-boot key held in a CPU register that is never spilled
to memory (a BEAM scheduler's reserved XMM register is a natural choice),
and $n$ is the integrity-bit budget --- 24 bits for a 36-bit address
space, 28 bits if we accept a 32-bit address cap.

The MAC is computed at every site that creates a term-bearing pointer:
allocation of a new boxed term, GC copy of a live term to a new arena,
construction of a new tuple element. It is verified at every site that
dereferences such a pointer.

A mismatch indicates that the bytes at the pointed-to address differ
from the bytes that were there when the pointer was created --- which
is exactly the signature of a RAM replay against an interior heap
location.

\subsection{Bit budget}

For a 64-bit machine word with BEAM's canonical 4-bit major tag:

\begin{center}
\begin{tabular}{lr}
\toprule
\textbf{Field} & \textbf{Bits} \\
\midrule
Major term tag (immediate / list / boxed / \dots) & 4 \\
Address bits (36-bit physical address space, $\le$~64~GB) & 36 \\
\textbf{Integrity tag (MAC)} & \textbf{24} \\
\midrule
Total & 64 \\
\bottomrule
\end{tabular}
\end{center}

A 24-bit MAC is short by cryptographic standards. The relevant security
question is not "can the attacker produce a forgery given an oracle"
but "can the attacker, in the time window between memory capture and
memory replay, find a stale ciphertext whose 24-bit MAC happens to
match the new MAC stored in the live pointer".  With a per-boot key
and content that changes at every minor GC ($\le$~100~ms), an attacker
must succeed within that window. The forgery probability per attempt
is $2^{-24}$; against a node serving even 100 requests per second, the
expected number of attempts before a successful blind forgery is
$\sim$~$10^7$, which translates to many minutes of focused attack on a
moving target. In practice the attacker is also constrained by their
DRAM-tap throughput and by the necessity of guessing which captured
ciphertext to replay; the effective security is materially higher than
the bit count alone suggests.

\subsection{MAC algorithm}

The choice of MAC is a performance--security trade-off, not a cryptographic
freshness question (the freshness comes from the per-boot key plus the
content-hash structure, not from MAC strength). Candidates:

\begin{itemize}[label=---]
\item \textbf{CRC-32C}: $\sim$~1 cycle per cache line on modern CPUs
with the SSE 4.2 \texttt{crc32} instruction. Not a MAC: an attacker who
can profile finds collisions deterministically. \emph{Insufficient.}
\item \textbf{HMAC-SHA-256}: cryptographically strong but $\sim$~hundreds
of cycles per evaluation. \emph{Unaffordable on the dereference hot path.}
\item \textbf{Truncated AES-128 PRF}: one AES round with AES-NI is
$\sim$~4 cycles; a full encryption is $\sim$~25--40 cycles. Truncated
to 24 bits, this composes as a near-ideal random function under standard
PRF assumptions. \emph{Likely the right choice.}
\item \textbf{SipHash-2-4}: $\sim$~10--20 cycles, designed for short
keyed hashes, software-only. Reasonable fallback if AES-NI is for
some reason unavailable.
\end{itemize}

The per-boot key $K$ is derived at startup from a high-entropy source
(\texttt{getrandom(2)} or \texttt{RDSEED}), placed in an XMM register
reserved by the BEAM scheduler dispatch loop, and never written to
memory. A check-failure trap clears the register before any debug
output is emitted.

\section{The graph-rooting argument}

The defence's elegance comes from the structure of BEAM's term graph.
Every term-bearing pointer is reachable from a small set of roots:

\begin{itemize}[label=---]
\item \textbf{Process control blocks (PCBs)}, one per Erlang process.
A PCB holds the root pointer to that process's heap, mailbox queue,
register save state, and scheduler queues.
\item \textbf{ETS table headers}, one per ETS table. The header holds
the root of the table's internal trees / hash tables.
\item \textbf{The atom table}, a single global structure mapping atom
identifiers to atom-name strings.
\item \textbf{Module code area}, a static area holding loaded BEAM
bytecode.
\item \textbf{Registered process / port / monitor / link tables},
small global tables.
\end{itemize}

If the attacker replays an \emph{interior} term --- a tuple element, a
list cell, a refc-binary header --- the parent pointer's MAC was
computed against the term's pre-replay bytes. The next dereference of
that parent pointer recomputes the MAC against the post-replay bytes
and detects the mismatch. Recursively, the attacker would have to
replay every ancestor up to the root; but each ancestor is itself
reached through a parent pointer carrying a MAC of its own pre-replay
bytes. The cascade terminates only at a root that has no parent
pointer.

This means the attack surface for \emph{undetectable} replay collapses
from "every byte of the BEAM heap" to "the root tables listed above".
The roots are small (PCB $\approx$~512 bytes; ETS table header
$\approx$~256 bytes; atom table $\approx$~tens of KB), bounded in count,
and accessed frequently by the scheduler.

\section{Keeping roots in cache}

The graph-rooting analysis still permits replay against any of the
root structures themselves, because they are not pointed at by a
parent. To close that gap we need the roots' bytes to never reach
DRAM in the first place.

The proposed mechanism is to keep the root structures resident in
on-die cache (L1 or L2) for the lifetime of the BEAM scheduler that
owns them, so that the attacker's bus tap on the DRAM channel never
observes the ciphertext. Linux does not expose direct cache pinning
to userspace, but two practical mechanisms approximate it:

\begin{itemize}[label=---]
\item \textbf{Access-pattern pinning.} The scheduler dispatch loop
already touches its current PCB on every reduction tick; ensuring that
the \emph{root pointer} fields of all schedulable PCBs are touched at
the same cadence keeps them in the L1/L2 LRU set. A small
``hot-PCB-header'' restructure --- placing the root pointer + integrity
key into a single cache line per PCB, and prefetching that cache line
on every scheduler-loop iteration --- makes this discipline cheap and
explicit.
\item \textbf{Intel CAT (Cache Allocation Technology).} Where available
(Xeon-class silicon mostly), CAT supports reserving a portion of the
last-level cache for a specific process. Less applicable to commodity
laptop hardware, but worth noting.
\end{itemize}

With root structures held in cache and interior structures protected by
tagged-pointer MACs, the residual attack surface is the small set of
\emph{cold} root structures that fall out of cache under load (e.g. PCBs
of dormant processes). For these, a periodic re-MAC-and-verify pass
against a separate per-table MAC fingerprint is feasible at coarse
granularity (tens of milliseconds), which is sufficient given the
attacker's tap-throughput constraints.

\section{Engineering considerations}

\subsection{Performance}

Every term dereference adds: a mask-off of the high bits, a MAC
recomputation against the dereferenced bytes, a comparison, a
conditional branch. On the BeamAsm JIT path with AES-NI, this can be
implemented as four instructions plus a fault on mismatch; on the
interpreter path it is more invasive. Ballpark estimates from
analogous CHERI / MTE work suggest a 10--30\% throughput cost is
realistic for a careful implementation, with the high end seen on
list-traversal-heavy benchmarks. A naïve port to the interpreter
without JIT specialisation can degrade by 2$\times$. Real numbers
require measurement on a prototype.

\subsection{Garbage collection}

BEAM's generational copying GC moves live terms during minor and major
collections. Each move invalidates the integrity tag of every parent
pointer to the moved object, because the bytes are now at a different
address (and, depending on the MAC variant, the MAC may also depend on
the address). Two design choices:

\begin{itemize}[label=---]
\item \emph{MAC over content bytes only.} MACs survive GC moves; the
GC need not recompute them. Cost: equal-content terms are
indistinguishable, so an attacker can replay an old \texttt{\{ok, 42\}}
over a new \texttt{\{ok, 42\}}, and the MAC matches. For
many-equal-shaped-terms workloads (counters, state machines) this is
a genuine residual surface.
\item \emph{MAC over content + address.} MACs are GC-recomputed but the
attacker cannot replay even shape-identical terms. Cost: GC work
increases significantly, especially for major collections that move
many objects.
\end{itemize}

A reasonable hybrid is content-only MACs for immutable / shared terms
(atoms, refc binaries, code) and content+address MACs for process-local
mutable terms. This splits the GC cost from the security benefit
appropriately, but it adds protocol complexity that needs careful
validation.

\subsection{Reaction policy}

When a MAC mismatch is detected, the system has options:

\begin{itemize}[label=---]
\item \textbf{Immediate halt}, kernel panic with an attestation-signed
fault note.
\item \textbf{Reboot}, with the next attestation envelope reporting the
fault digest in PCR 15.
\item \textbf{Crash the affected Erlang process only}, supervised
restart, log the fault in the runtime event log so that the next
attestation envelope carries it. This option is especially attractive
for LapEE because the fault digest itself becomes a cryptographic
record of the attack: a verifier sees an envelope claiming
\texttt{memory-replay-detected} with a hash binding the fault to the
attesting node's identity. That is a much more useful signal than a
silent reboot.
\end{itemize}

The third option preserves the attestation channel through the attack,
which is exactly what a remote verifier wants.

\subsection{Compatibility with vanilla OTP}

The patch changes BEAM's pointer representation; modules compiled
against vanilla BEAM and loaded into a hardened runtime must either
be recompiled (the safer route) or the runtime must transparently
handle ``untagged'' pointers via a compatibility mode. We recommend
the recompile route for LapEE-only deployment, accepting that a
hardened LapEE node runs only LapEE-compiled BEAM bytecode.

\subsection{Estimated cost of a production-quality patch}

Two to four engineer-months for the core BEAM modifications, plus a
comparable amount for measurement, fuzzing, and OTP regression
testing. A research prototype demonstrating detection on synthetic
replay against process heaps (no GC integration, no JIT
specialisation, interpreter-only) is on the order of two to three
weeks for someone fluent in BEAM internals.

\section{Open design questions}

\begin{enumerate}
\item \emph{Address layout.} What address-bit cap is acceptable?
36 bits gives 64~GB and 24-bit MAC; 32 bits gives 4~GB and 28-bit MAC.
The 32-bit cap is restrictive for general use but acceptable for
LapEE laptops if we declare 4~GB enough for an attestation appliance.
\item \emph{MAC scope.} Per-pointer MAC of pointed-to-bytes-only, or
of bytes-plus-address, or a hybrid as above? This requires
benchmarking against representative HyperBEAM workloads (HTTP request
handling, attestation, ETS-heavy code).
\item \emph{Atom table protection.} The atom table is a single global
structure that grows monotonically. Protecting it with the same
per-pointer scheme is overkill; a periodic Merkle root computed over
the atom-name-to-id map and held in a register would be more
efficient.
\item \emph{Code area protection.} Loaded BEAM code is essentially
read-only after load. A simple per-module MAC computed at load time
and verified periodically is sufficient; this is independent of the
per-pointer scheme.
\item \emph{Multi-scheduler coordination.} Each scheduler's reserved
register holds the per-boot key. Are there cross-scheduler dereferences
(e.g. reading a remote PCB) that need a shared key view, and how is
that shared without leaking $K$ to memory?
\end{enumerate}

\section{Path to a prototype}

We propose the following sequence:

\begin{enumerate}
\item \textbf{Vendored OTP fork.} Branch the OTP release matching the
deployed HyperBEAM (\texttt{27.3.4}). Add a build-time configure flag
\texttt{--enable-tagged-mac-replay-defence}.
\item \textbf{Minimum viable scope.} Implement per-pointer MAC for
boxed terms in the process heap only. No atom-table, ETS, or code-area
coverage in the prototype. Use truncated AES-128 with the per-boot key
in an XMM register.
\item \textbf{Single dereference path.} Modify \texttt{HALLOC} and the
boxed-pointer-deref macros in the interpreter; skip BeamAsm. This
gives an honest baseline without micro-optimisation.
\item \textbf{Synthetic replay harness.} Build a test that allocates
known terms, captures their bytes, then writes the captured bytes back
to the same heap address while a parent process attempts to dereference.
Verify the mismatch is detected.
\item \textbf{Measurement.} Run the OTP benchmark suite plus a
representative HyperBEAM workload (\texttt{hb\_http} request flow with
\texttt{dev\_tpm2} \texttt{/attestation} calls) and report throughput,
GC pause, and tail latency vs.\ vanilla OTP.
\item \textbf{Decision point.} On the basis of measured numbers, decide
whether to proceed with full BeamAsm specialisation, GC integration,
and root-structure caching, or to publish the prototype as a research
note and wait for hardware integrity to become available on the
target platform.
\end{enumerate}

\section{Conclusion}

The construction is not novel in its ingredients --- tagged-pointer
integrity has appeared in CHERI, ARM MTE, SoftBound, and several
research systems --- but its application to BEAM specifically, in
service of a software-only defence against TME-replay attacks on
commodity laptop hardware, appears to be unaddressed in the literature.
The graph-rooting analysis is what makes the construction interesting:
the attack surface reduces to a small bounded set of roots, and
keeping those roots cache-resident closes the residual gap to a
practically infeasible target. The performance cost is real but
likely acceptable for a node whose primary workload is HTTP request
handling at ordinary rates rather than CPU-bound computation.

For LapEE today this is future work. The current build hardens what
it can with built-in address-space and allocator randomisation; the
tagged-MAC patch is documented here as the proper next-generation
defence and as a candidate research contribution in its own right.

\end{document}
